% 00_abstract.tex — 摘要

本研究系统分析了非 LEO 地球轨道（MEO、GEO、HEO）以及地日/地月拉格朗日点（L1、L2、L4、L5）部署太空算力基础设施的可行性，涵盖辐射环境、热环境、\deltav 发射成本、太阳能供电、通信延迟、在轨维护和全生命周期成本（TCO）七个维度。

研究发现，\textbf{LEO 并非太空算力的最优轨道}。GEO（地球同步轨道）在综合成本指数评估中以 0.83 的得分（LEO 为 1.00，数值越低越优）排名第一，其 15 年已验证的硬件寿命、即将可用的在轨服务能力以及简化的热设计和光伏需求，使其在"可维护性"和"全生命周期"维度上对 LEO 具有根本性优势。地月 L4/L5 凭借三体引力稳定性和零站位保持需求，是最具长期战略价值的深空算力平台。地日 L2 的 40 K 天然冷端为未来量子-经典协同计算提供了独特机遇。

研究同时揭示了"\tid 悖论"：地日 L1/L2 在适度屏蔽后的年 \tid 仅为 2--5~\kradyr，低于穿越南大西洋异常区（SAA）的 LEO 轨道（10--30~\kradyr），但 \seu 率比 LEO 高 100--1000 倍。这意味着深空算力的辐射挑战本质是 \seu 挑战而非 \tid 挑战。MEO 和 HEO 由于极端辐射环境和缺乏独特应用优势，不推荐作为独立太空算力部署轨道。

非 LEO 太空算力最可能的演进路径为：GEO 先行（2028--2035 年）\textrightarrow\ 地月 L4/L5 跟随（2032--2040 年）\textrightarrow\ 地日 L2 利基应用（2035 年以后）。
